% Source: tex/chapters/ch05/08_テストプロセス.tex
\subsubsection{テストサブプロセスと活動}
\paragraph{テスト計画プロセス}
テスト計画では、担当者の調整、テスト目的、完了基準、設備、文書、リスク計画を定義する。組織方針、テスト段階、テスト種類、テスト環境、実行、監視、完了、報告を整える。

\paragraph{テスト設計と実装}
テスト設計と実装では、テストレベルと選んだ技法に基づいてテストケースを作る。ツールを使う場合、ソースコード、仕様、データ構造、テスト基準などを入力し、テストスイートを生成する。場合によっては、期待結果もオラクル機能で生成できる。

\paragraph{テスト環境の構築と保守}
テスト環境は、テストを実行するための設備、ハードウェア、ソフトウェア、ファームウェア、手順を含む。シミュレーション環境、本番に近い環境、監視・ログ機能などを整え、他のソフトウェア工学ツールと互換性を持たせる必要がある。

\paragraph{統制された実験とテスト実行}
テスト実行は、科学的な統制実験の原則を意識すべきである。別の人が結果を再現できるほど、手順、環境、SUT の版、入力、観測結果を明確に記録する。受入テストでは、A/B テストのように利用者の選好を統計的に評価することもある。

\paragraph{テストインシデント報告}
テストインシデント報告は、いつ、誰が、どの構成でテストを行い、どの結果が得られたかを記録する。予期しない結果は問題報告システムに記録され、後のデバッグや変更管理の入力になる。すべての異常が欠陥とは限らないため、分析により原因を特定する必要がある。
